\section{Technology Stack Overview}

\subsection{Technology Selection Rationale}
\textbf{Java} was selected for backend and plugin development due to the maturity and stability of the JVM ecosystem, as well as its widespread adoption. Additionally, the Minecraft server ecosystem and its plugin frameworks are primarily designed for Java, making it the most natural and well-supported choice for extending gameplay functionality.

Web technologies were chosen to enable a modern and interactive user interface. The dashboard is built using \textbf{React, TypeScript, and Tailwind CSS}, leveraging their strong community support, market relevance, and suitability for building scalable and maintainable front-end applications. These technologies also provide an opportunity to apply and deepen practical knowledge of widely adopted industry tools.

\textbf{Docker} was adopted to ensure consistency across development and deployment environments. By containerizing each module, the project simplifies setup, reduces environment-related issues, and enables reliable and reproducible deployments.

\subsection{Backend Technologies}
\textbf{Java 17 \& Spring Boot 3.5.7}: Development of RESTful APIs and dependency management.

\textbf{PostgreSQL 18}: Relational database for persistent data storage.

\textbf{Flyway}: Database schema versioning and migration management.

\subsection{Frontend Technologies}
\textbf{React}: Component-based user interface development.

\textbf{JavaScript \& TypeScript}: Improved type safety, maintainability, and developer experience.

\textbf{Tailwind CSS}: Utility-first framework for fast and consistent styling.

\subsection{Infrastructure and DevOps Technologies}
\textbf{Gradle}: Build automation and dependency management for Java projects.

\textbf{Docker 29}: Application containerization and environment standardization.

\subsection{Versioning and Compatibility Strategy}
To be defined.
